\chapter[Careless Responding Detection]{Careless Responding Detection: The Field, Current Practices, and Software}
\label{ch:intro}

\section{The Problem of Careless Responding}

Self-report questionnaires rest on a tacit assumption: that respondents read
each item, retrieve the relevant information, form a judgment, and map it onto
the response scale. When this process breaks down---when respondents answer
without attending to item content---the resulting data carry the surface form
of valid measurement but none of its substance. This phenomenon, variously
labeled \emph{careless responding}, \emph{insufficient effort responding}
(IER), \emph{inattentive responding}, or \emph{random responding}, has become
one of the most actively studied threats to the validity of survey research
\parencite{meade2012identifying, curran2016methods, ward2023dealing}.

The terminological proliferation reflects subtly different emphases.
\textcite{huang2012detecting} introduced \emph{insufficient effort responding}
to foreground the motivational mechanism: respondents who are unwilling to
expend the cognitive effort the task demands. \textcite{krosnick1991response}
had earlier framed the same behavior as \emph{satisficing}---respondents who,
rather than \emph{optimizing} their answer through the full cognitive sequence,
settle for a response that is merely good enough, or in the limiting case of
\emph{strong} satisficing, select responses with little or no reference to item
content. \emph{Careless responding} is the broader umbrella term now favored in
methodological reviews \parencite{ward2023dealing}; it encompasses both the
low-effort respondent and the respondent whose attention lapses intermittently
over a long instrument. Throughout this thesis I use \emph{careless responding}
as the cover term and reserve more specific labels for particular response
patterns.

Whatever the label, the empirical consequences are well documented and
consistently damaging. Careless responding inflates measurement error and
depresses scale reliability \parencite{huang2012detecting}. Because careless
responses are, to a first approximation, unrelated to the constructs a
questionnaire targets, they attenuate true relationships between variables; yet
they can also \emph{create} spurious relationships, most insidiously when
careless respondents endorse the midpoint or an extreme of every scale and
thereby manufacture artifactual covariance across otherwise unrelated measures
\parencite{huang2015insufficient}. \textcite{maniaci2014caring} demonstrated
across several studies that even modest proportions of inattentive participants
can meaningfully distort effect sizes, suppress genuine effects, and reduce
statistical power. \textcite{kam2015careless} showed that careless and
acquiescent responding can distort the very factor structure recovered from a
set of items, splitting a unidimensional construct into spurious method
factors. In clinical and applied settings the stakes are higher still:
\textcite{conijn2019assessment} found that careless responding in routine
outcome monitoring biased symptom estimates in ways that could affect treatment
decisions.

Estimates of prevalence vary with population, mode of administration, and
detection method, but they are rarely negligible. Reviews place the typical
incidence somewhere between roughly 5\% and 15\% of respondents in
well-controlled samples, with substantially higher rates in unsupervised online
panels and crowdsourced labor markets \parencite{ward2023dealing,
meade2012identifying, curran2016methods}. \textcite{bowling2016who} added an
important nuance: carelessness is not purely situational noise but is partly a
stable individual difference, correlated with low conscientiousness and low
agreeableness, meaning that the same respondents tend to provide low-effort data
across studies. The practical upshot, emphasized in every recent treatment of
the topic, is that careless responding cannot be assumed away: it must be
actively prevented where possible and detected and addressed where not
\parencite{desimone2015best, desimone2018dirty, ward2023dealing}.

\section{A Taxonomy of Causes and Response Patterns}

Effective detection requires a clear picture of \emph{what careless responding
looks like} in the data, because different careless behaviors leave different
statistical fingerprints. The literature converges on a useful organizing
distinction between \emph{consistent} and \emph{inconsistent} careless
responding \parencite{ward2023dealing, curran2016methods}.

\emph{Consistent} careless patterns are those in which the respondent produces
a regular, low-information sequence. The clearest example is
\emph{straightlining} (also called \emph{non-differentiation}): selecting the
same response option for a long run of consecutive items, often the scale
midpoint or an extreme. A closely related notion is the \emph{longstring}---the
maximum length of consecutive identical responses, which \textcite{johnson2005ascertaining}
established as a simple but effective marker of inattentive protocols.
\emph{Acquiescent} responding (uniform agreement regardless of item keying) is
another consistent pattern, particularly damaging because it survives naive
reliability checks while corrupting the covariance structure
\parencite{kam2015careless}. Consistent patterns are, paradoxically, both easy
and hard to detect: easy because their regularity is conspicuous to a longstring
counter, hard because a respondent who endorses the midpoint of every item can
look indistinguishable from a genuinely neutral one.

A methodological caveat sharpens this picture and bears directly on the design
of the detector developed in this thesis. Best practice in questionnaire
construction is to \emph{randomize the order in which items are presented} to
each respondent, which controls for order and context effects and breaks up
contiguous blocks of same-construct items \parencite{ward2023dealing}. This
otherwise desirable practice substantially undermines the longstring index. A
careless respondent who repeatedly selects the same option no longer produces a
long \emph{contiguous} run of identical responses once the data are returned to
their canonical, construct-grouped order for analysis: the repeated answers are
scattered across non-adjacent positions, and reverse-keyed items, after
recoding, break whatever run remains. Under randomized administration the
longstring counter therefore loses much of its sensitivity, and even
straightlining---the very pattern it was designed to catch---can slip past it.
This is one reason the simulation studies in later chapters model longstrings as
\emph{non-contiguous}, scattered across the response vector rather than as neat
adjacent blocks: a deliberately conservative choice that reflects modern
randomized administration and that motivates detectors whose signal does not
depend on the contiguity or ordering of responses.

\emph{Inconsistent} careless patterns, by contrast, involve responses that vary
across items but bear no coherent relationship to item content. The prototype is
\emph{purely random} responding, in which answers are effectively drawn at
random from the response scale. Such protocols display normal-looking variance
and no long runs, so they slip past longstring and variability checks; what
betrays them is the \emph{absence of the correlational structure} that attentive
respondents exhibit---answers to items measuring the same latent trait, which
should covary, instead scatter \parencite{meade2012identifying, curran2016methods}.
\emph{Mixed} patterns, in which a respondent attends to part of the instrument
and disengages for the rest (for example through mid-survey
\emph{fatigue}), combine features of both families and are correspondingly
difficult to catch with any single index \parencite{ward2023dealing}.

This taxonomy matters for a reason that runs through the entire thesis: the
indices in widest use were largely designed for, and excel at, consistent
careless patterns, whereas inconsistent patterns---random and mixed
responding---are precisely where most methods are weakest. The cross-factor
coherence approach developed in later chapters is motivated by this gap.

\section{Prevention by Design}

Because no detection method is perfect, the first line of defense is to design
instruments and data-collection procedures that discourage careless responding
and that embed signals of inattention directly into the questionnaire
\parencite{ward2023dealing, desimone2015best}.

The most common device is the \emph{attention check}, which takes several forms.
\emph{Instructed-response items} (IRIs) direct the respondent to select a
specific option (``Please select \emph{strongly agree} for this item''); a
deviation flags inattention. \emph{Bogus items} assert something patently true
or false (``I was born on February 30th''); an implausible answer is
diagnostic. \emph{Instructional manipulation checks} (IMCs), introduced by
\textcite{oppenheimer2009instructional}, embed a comprehension test within a
seemingly ordinary item and have been shown to identify respondents who are
satisficing rather than reading. A self-report single-item approach simply asks
respondents, at the end, whether their data should be used---a surprisingly
informative question when posed without penalty \parencite{meade2012identifying}.

These instruments are valuable but limited, and the limitations are
instructive. \textcite{curran2019paid} showed that attention-check items
generate a non-trivial rate of \emph{valid-but-incorrect} responses: attentive
participants sometimes fail them (through misreading or genuine ambiguity),
while determined careless respondents can pass an isolated check by chance or by
briefly attending to a conspicuous item. A single check therefore has limited
sensitivity and imperfect specificity. Attention checks also consume
questionnaire space, can alter respondent behavior by signaling surveillance,
and---being placed at fixed positions---catch inattention only at those points,
missing the intermittent disengagement characteristic of fatigue. For
time-stamped or computer-administered surveys, \emph{response latency} offers a
complementary design-based signal: implausibly fast page completions indicate
that a respondent could not have read the items, a logic formalized by
\textcite{wise2005response} in the \emph{response time effort} measure and later
embedded in model-based approaches \parencite{ulitzsch2022response}. Response
times, however, are unavailable in the many settings where only the final
responses are recorded.

The recurring lesson is that prevention reduces but does not eliminate careless
responding, and that the design-based signals it provides are partial. This
motivates \emph{post-hoc} detection from the response data themselves---methods
that require no special items, no timing information, and no advance planning,
and that can therefore be applied to data already collected.

\section{Post-hoc Identification Methods}

Post-hoc methods infer carelessness from the structure of the observed
responses. They fall into several families, each exploiting a different
statistical signature of inattention. The two most authoritative syntheses of
this literature---\textcite{meade2012identifying} and \textcite{curran2016methods}---organize
the field around a core set of indices that I summarize here and collect, with
their software implementations, in Table~\ref{tab:indices}.

\subsection{Consistency and Response-Pattern Indices}

\emph{Longstring} counts the maximum number of consecutive identical responses
in a protocol. It is the canonical detector of straightlining and is trivial to
compute, but it is blind to inconsistent careless responding and is biased by
the number of items keyed in the same direction \parencite{johnson2005ascertaining,
meade2012identifying}.

The \emph{intra-individual response variability} (IRV) is the standard deviation
of a respondent's answers across a set of items \parencite{dunn2018intra}. Very
low IRV signals straightlining (little or no variation), while implausibly high
IRV can indicate random responding across the full scale range. Because it
condenses a protocol to a single dispersion statistic, IRV is cheap and
intuitive, but it discards all information about \emph{which} items were answered
how.

\emph{Psychometric synonyms and antonyms} \parencite{meade2012identifying}
operationalize coherence at the item-pair level. One identifies pairs of items
that are strongly positively correlated in the sample (synonyms) or strongly
negatively correlated (antonyms), then computes, \emph{within each respondent},
the correlation between their answers to the two members of each pair. An
attentive respondent answers synonym pairs similarly and antonym pairs
oppositely, producing a high within-person correlation; a careless respondent
does not. This is a powerful idea because it targets inconsistent careless
responding directly, and it is the conceptual ancestor of the cross-factor
coherence method developed later in this thesis.

The \emph{person--total correlation} correlates a respondent's profile of
answers with the sample's mean response to each item; respondents who answer
against the grain of the typical pattern receive low values. \emph{Even--odd
consistency} splits a scale into even- and odd-numbered items, computes two
subscale scores, and assesses their agreement; \emph{resampled individual
reliability} generalizes this by repeatedly splitting and correlating. These
consistency indices trace back to classical work on individual-protocol
validity \parencite{jackson1976jackson, johnson2005ascertaining}.

\subsection{Multivariate Outlier Detection}

A complementary strategy treats careless protocols as multivariate
\emph{outliers}. The \emph{Mahalanobis distance} \parencite{mahalanobis1936generalised}
measures how far a respondent's response vector lies from the sample centroid,
scaled by the inverse covariance matrix, and is the most widely used outlier
index in this literature \parencite{meade2012identifying, curran2016methods}.
Under multivariate normality the squared distance follows a chi-square
distribution with degrees of freedom equal to the number of items, which
suggests a principled cutoff. In practice, Likert data violate the normality
assumption badly, so the chi-square reference distribution is only approximate,
and---as the simulation chapters of this thesis document in detail---the
nominal cutoff can be far too conservative for questionnaire data, flagging
almost no one. Mahalanobis distance also detects \emph{any} multivariate
outlier, careless or not, including genuinely unusual but attentive respondents.

\subsection{Person-Fit and Item-Response-Theory Approaches}

A distinct tradition grew out of item response theory (IRT), where
\emph{person-fit} statistics quantify how well an individual's response pattern
conforms to the fitted measurement model. The $l_z$ statistic and the Rasch
\emph{infit}/\emph{outfit} indices are the best known; \textcite{karabatsos2003comparing}
compared thirty-six such statistics for detecting aberrant response patterns,
and \textcite{meijer2016practical} provided a practical guide to applying them
in clinical research. Person-fit methods rest on a strong measurement model,
which is their strength (a principled probabilistic account of ``aberrance'')
and their weakness (sensitivity to model misspecification, and a need for the
model to be estimable). More recently, \emph{mixture-model} formulations treat
the population as a blend of attentive and careless latent classes and estimate
class membership directly; \textcite{ulitzsch2022response} combined item
responses with response times in a latent-response mixture model that
simultaneously identifies and adjusts for careless and insufficient-effort
responding.

\subsection{Machine-Learning and Ensemble Approaches}

The newest family abandons the idea of a single hand-crafted index in favor of
\emph{learning} a careless/attentive boundary from many indices at once.
\textcite{schroeders2022detecting} trained a stochastic gradient-boosting
classifier on a battery of careless indicators and reported detection
performance exceeding that of any constituent index, demonstrating that the
indices carry \emph{complementary} information that a flexible classifier can
exploit. This ensemble logic is directly relevant to the present thesis, whose
later chapters combine a novel coherence index with established detectors inside
a Random Forest \parencite{breiman2001random}. The turn toward machine learning
has not gone uncontested. \textcite{alfons2024openscience} offer an
open-science critique, warning that supervised approaches risk overfitting to
the idiosyncrasies of the simulated or instructed-careless data on which they
are trained, that ground-truth labels for careless responding are inherently
uncertain, and that reproducible benchmarks and shared data are largely absent
from the field. Their related work on identifying the \emph{onset} of careless
responding within a protocol \parencite{welz2023onset} illustrates an
alternative, model-based direction. These cautions frame the methodological
commitments of this thesis---in particular the emphasis on validation across
independent simulation draws and on deployment-realistic, label-free
calibration.

\section{Cutoffs, Thresholds, and Combining Indices}

Computing an index is only half the task; the other half is deciding \emph{where
to draw the line}. Every method above yields a continuous score, and converting
that score into a binary careless/attentive flag requires a threshold. This
seemingly technical step is in fact one of the least settled and most
consequential issues in the field.

Two broad strategies exist. The first uses a \emph{theoretical} cutoff derived
from a reference distribution---most prominently the chi-square cutoff for
Mahalanobis distance---which has the appeal of requiring no tuning but, as
noted, rests on distributional assumptions that questionnaire data routinely
violate. The second uses an \emph{empirical} cutoff: a fixed quantile of the
observed scores, a visual ``elbow'' in the score distribution, or a value
optimized against known labels. Empirical cutoffs adapt to the data but risk
either circularity (when tuned on the same labels used to evaluate them) or
arbitrariness (when chosen by eye). \textcite{ward2023dealing} note that the
absence of consensus cutoffs is a genuine obstacle to cumulative science,
because two analysts applying ``the same'' index to the same data can reach
different conclusions about who is careless.

The comparative literature offers two robust findings that shape best practice.
First, no single index dominates across conditions:
\textcite{niessen2016detecting} compared several methods on web-questionnaire
data and concluded that the best choice depends on the type of careless
responding present and on questionnaire characteristics, with different indices
catching different careless respondents. Second---and following directly from
the first---\emph{combining} indices outperforms relying on any one of them
\parencite{curran2016methods, schroeders2022detecting}. The recommended
practice in recent reviews is therefore to compute multiple indices, examine
their convergence, and flag respondents identified by several of them
\parencite{ward2023dealing, desimone2015best}. The ensemble and threshold-calibration
methods developed in this thesis are a formalization of exactly this advice: they
replace ad hoc multi-index agreement rules with a trained classifier and a
deployment-realistic operating point.

\section{Software for Careless Responding Detection}

The methods reviewed above are only as useful as the tools that implement them,
and here the \textsf{R} ecosystem has become the de facto standard for careless
responding detection. This section surveys the packages in current use; their
indices and the constructs they target are summarized in
Table~\ref{tab:indices}.

The central package is \textsf{careless} \parencite{yentes2023careless}, which
provides a unified implementation of the core response-pattern indices:
\texttt{longstring} (and its average-run variant), \texttt{irv}, \texttt{mahad}
(Mahalanobis distance, with optional robust covariance estimation),
\texttt{evenodd} (even--odd consistency), and \texttt{psychsyn}/\texttt{psychant}
(psychometric synonyms and antonyms). Because it gathers the indices of
\textcite{meade2012identifying} and \textcite{curran2016methods} behind a
consistent interface, \textsf{careless} is the package most often cited in
applied work and the natural baseline against which new methods are compared.

Several other packages contribute pieces of the workflow. The
\textsf{responsePatterns} package \parencite{rihacek2023responsepatterns}
implements autocorrelation-based screening for repetitive response patterns
\parencite{gottfried2022autocorrelation}, targeting the machine-like regularity
that some careless protocols exhibit. For the multivariate-outlier branch, the
\textsf{performance} package \parencite{ludecke2021performance} exposes
\texttt{check\_outliers}, which computes Mahalanobis (and several robust
alternatives) within a tidy model-checking framework, while the venerable
\textsf{psych} package \parencite{revelle2024psych} supplies the
factor-analytic and reliability machinery---exploratory factor analysis,
parallel analysis, and scale-reliability estimation---on which several careless
indices and most downstream psychometric analyses depend. Mixture-model
approaches to careless detection draw on general-purpose finite-mixture engines
such as \textsf{mclust} \parencite{scrucca2016mclust} and \textsf{mixtools}
\parencite{benaglia2009mixtools}, which fit the latent attentive/careless class
structure used by model-based methods. All of these run within base
\textsf{R} \parencite{rcoreteam2024}, and supervised-ensemble approaches add a
classifier layer---gradient boosting in \textcite{schroeders2022detecting}, or a
Random Forest in the present work---on top of the index computations these
packages provide.

What the ecosystem still lacks is as informative as what it offers.
\textcite{alfons2024openscience} observe that, despite the maturity of the
individual-index tooling, the field has no shared, labeled benchmark dataset
against which competing detectors can be evaluated on equal terms, and that
machine-learning detectors in particular are typically validated only on the
simulated data their authors generate. There is, similarly, no standard,
well-validated routine for \emph{calibrating} an index threshold without access
to ground-truth labels---the very problem the deployment-oriented chapters of
this thesis address. Outside \textsf{R}, careless detection is comparatively
underserved: commercial survey platforms offer little beyond fixed-position
attention checks, and although general-purpose machine-learning libraries in
other languages could in principle implement any of these methods, no
established, careless-specific toolkit exists outside the \textsf{R} packages
surveyed here.

\begin{table}[t]
\centering
\caption{Principal careless-responding detection methods, the careless
behavior each is best suited to detect, and their implementation in current
\textsf{R} software.}
\label{tab:indices}
\footnotesize
\begin{tabular}{L{0.20\linewidth} L{0.30\linewidth} L{0.14\linewidth} L{0.25\linewidth}}
\toprule
\textbf{Method} & \textbf{What it captures} & \textbf{Target pattern} & \textbf{R implementation} \\
\midrule
Longstring & Longest run of identical consecutive responses & Consistent (straightlining) & \texttt{careless::\allowbreak longstring} \\
IRV & SD of a respondent's answers across items & Consistent (low) / random (high) & \texttt{careless::\allowbreak irv} \\
Psychometric synonyms / antonyms & Within-person correlation over high-$|r|$ item pairs & Inconsistent (random, mixed) & \texttt{careless::\allowbreak psychsyn},\newline \texttt{psychant} \\
Person--total correlation & Agreement of a profile with the sample mean profile & Inconsistent & \texttt{careless} / base \textsf{R} \\
Even--odd / resampled reliability & Internal consistency across item splits & Inconsistent & \texttt{careless::\allowbreak evenodd} \\
Mahalanobis $D^2$ & Multivariate distance from the sample centroid & Multivariate outliers & \texttt{careless::\allowbreak mahad};\newline \texttt{performance::\allowbreak check\_outliers} \\
Autocorrelation screening & Machine-like repetitiveness of a protocol & Consistent (patterned) & \texttt{responsePatterns} \\
Person-fit ($l_z$, infit/outfit) & Departure from a fitted IRT model & Aberrant patterns & \texttt{mirt} / \texttt{PerFit} \\
Mixture / response-time models & Latent attentive vs.\ careless classes & All (model-based) & \texttt{mclust}, \texttt{mixtools} \\
Ensemble / machine learning & Learned combination of many indices & All & gradient boosting; Random Forest \\
\bottomrule
\end{tabular}
\end{table}

\section{Gaps in the Field and the Contribution of This Thesis}

This review motivates the work that follows by making three gaps explicit.

First, the methods in widest use are strongest on \emph{consistent} careless
responding and weakest on the \emph{inconsistent} patterns---random and
mixed---that are arguably the hardest to prevent and the most damaging to
covariance-based analyses. Longstring and IRV detect regularity; Mahalanobis
distance detects outlyingness; but the cross-item \emph{coherence} that
distinguishes an attentive from a randomly responding participant is captured
only partially, by the synonym/antonym indices, and only for short lists of
hand-selected item pairs. A method that exploits the full correlational
structure of a multi-construct questionnaire to detect inconsistent careless
responding would fill a real gap. This is the motivation for the
permutation-based coherence index, \emph{ReReReRe}, developed in Chapter~3.

Second, the field has converged on the principle that \emph{combining} indices
beats any single one \parencite{niessen2016detecting, curran2016methods,
schroeders2022detecting}, but the dominant practice---flagging respondents
identified by several indices under separately chosen cutoffs---remains
informal. A trained classifier that learns the optimal combination of
complementary detectors, and that is evaluated honestly across independent
data, is the natural next step. This is the motivation for the Random Forest
ensemble of Chapter~4.

Third, every detection method ultimately requires a \emph{threshold}, and the
field lacks a validated procedure for setting one without ground-truth labels
that no analyst possesses in practice \parencite{ward2023dealing,
alfons2024openscience}. A calibration strategy that is both deployment-realistic
and competitive with label-dependent oracles would make these methods usable by
the applied researchers who need them. This is the motivation for the
rate-aware calibration evaluated in Chapter~5.

The remaining chapters take up these gaps in turn. Chapter~2 reviews the
statistical concepts on which the methods rest---correlation, exploratory factor
analysis, the Random Forest classifier, and the evaluation metrics used
throughout. Chapter~3 introduces ReReReRe and its underlying coherence
principle. Chapter~4 develops the ensemble that combines ReReReRe with the
established detectors surveyed above. Chapter~5 applies the ensemble under a
deployment-realistic protocol and characterizes when, and how well, it works.

